\documentclass[10pt,letterpaper]{moderncv}
\moderncvstyle{banking}
\moderncvcolor{black}
\definecolor{color1}{RGB}{0,0,0}% accent color
\definecolor{color2}{RGB}{0,0,0}% main color
\usepackage[margin=0.3in]{geometry}
\usepackage{enumitem}

\name{Jeff}{Calderon}
\address{Marysville}{WA}
\phone{916\,712\,8396}
\email{jeff.d.calderon@gmail.com}
\social[linkedin]{in/jeffdcalderon}
\social[github]{aggressor-FZX}
\nopagenumbers{}
\begin{document}
\makecvtitle

% TAGLINE / SUMMARY-------------------------------
\vspace*{-9mm}
\begin{center}
\textit{Engineering Data Scientist~|~Safety Analytics~|~Sensor Data~|~Predictive Modeling}
\end{center}
\vspace*{1mm}

Engineering data scientist with a physics foundation and 10+ years of applied experience designing predictive models, analyzing sensor and temporal data streams, and delivering actionable safety insights to operational and executive stakeholders. Background spans safety-critical federal research (Naval Research Laboratory national security programs), multi-sensor oceanographic data analysis at scale (10+ TB, NOAA research fleet), aviation-adjacent operations (U.S. Air Force aerospace ground equipment), and Boeing Defense-sponsored ML research. Experienced building ETL pipelines, statistical models, anomaly detection systems, and data visualization products that drive risk-based decisions. Physics degree with a peer-reviewed publication.

% TECHNICAL SKILLS-------------------------------
\vspace*{2mm}
\section{Technical Skills}
\vspace*{1mm}
\textbf{Statistical and Predictive Modeling:} Descriptive, diagnostic, predictive, and prescriptive model development, statistical analysis, Monte Carlo simulation, risk quantification, anomaly detection, time-series and trajectory analysis, algorithm design and performance evaluation, model validation, safety risk assessment, Bayesian concepts, regression, gradient boosting, random forest, neural networks

\textbf{Sensor and Temporal Data:} Multi-sensor data fusion, temporal data analysis, sensor stream ingestion and validation, large-scale time-series processing, ETL/ELT pipeline design, data quality and governance, geospatial data analysis (ArcGIS, CARIS), sonar data processing~-- directly analogous to ADS-B and FDM aviation sensor data workflows

\textbf{Languages and Tools:} Python (primary), SQL, C++, Bash, JavaScript, TypeScript, Pandas, NumPy, PyTorch, TensorFlow, Scikit-learn, Matplotlib, Seaborn, Power BI, Tableau, AWS (Lambda, Glue, Step Functions, S3), Azure, Docker, Git, Linux

\textbf{NLP and AI:} NLP (transformers, attention mechanisms, embeddings, sequence models), LLM integration (LangChain, OpenAI API, RAG), LLM fine-tuning, generative models (PyTorch DDP training on GPU clusters), NVIDIA H100/A100, SLURM HPC

% PROFESSIONAL EXPERIENCE-------------------------------
\vspace*{2mm}
\section{Professional Experience}
\vspace*{1mm}

\textbf{Data Analyst Intern~-- Washington State Data Exchange for Public Safety} \hfill 04/2025 -- Present
\vspace*{-2mm}
\begin{itemize}[label={-}, itemsep=-1.0ex]
  \item Design and implement ETL/ELT pipelines (AWS Step Functions, Glue, Lambda) ingesting multi-source safety and operational data; build Power BI dashboards and reports tracking safety-relevant KPIs for law enforcement and EMS leadership.
  \item Perform exploratory analysis and variance reporting on operational datasets; translate findings into actionable recommendations for agency directors and cross-functional stakeholders.
\end{itemize}

\textbf{Operations Manager~-- National Oceanic and Atmospheric Administration (NOAA)} \hfill 12/2019 -- 05/2023
\vspace*{-2mm}
\begin{itemize}[label={-}, itemsep=-1.0ex]
  \item Built descriptive and predictive analytical models synthesizing high-dimensional sensor data from 13 oceanographic research vessels; designed Python-based decision-support tools for operational safety and readiness monitoring, reducing manual reporting overhead by 40\%.
  \item Delivered safety-relevant data insights and risk-based recommendations to Regional Director and fleet leadership through regular executive briefings; translated complex multi-source analytical findings into clear, actionable management summaries.
  \item Launched Government Official Travel analytics portal: gathered requirements, designed data architecture, built monitoring system, and achieved 85\% adoption within the first month.
  \item U.S. Coast Guard certified Medical Person in Charge (MPIC) aboard NOAA research vessels: performed patient assessment, administered medications, and managed medical emergencies -- demonstrates applied safety risk assessment, regulatory compliance, and critical decision-making in high-stakes operational environments directly relevant to aerospace safety analytics.
\end{itemize}

\textbf{Marine Data Analyst~|~Hydrographer~-- NOAA Commissioned Officer Corps (Office of Coast Survey)} \hfill 12/2016 -- 05/2019
\vspace*{-2mm}
\begin{itemize}[label={-}, itemsep=-1.0ex]
  \item Processed and fused multi-sensor shipboard data streams in support of the U.S. Office of Coast Survey: multibeam sonar, single-beam sonar, side-scan sonar, LiDAR, GPS, inertial motion, tidal modeling, and shipboard weather sensor arrays -- a sensor fusion workload directly analogous to ADS-B and FDM aviation safety data pipelines.
  \item Designed and implemented anomaly detection algorithms on sonar and geospatial datasets to identify navigation hazards and seafloor irregularities for official nautical chart updates; evaluated algorithm performance against surveyor-validated ground truth and iterated based on results -- safety-critical algorithm development with real-world consequences if wrong.
  \item Applied trajectory-based analysis to vessel GPS and inertial motion data to reconstruct and validate survey track lines, assess data coverage, and flag positional anomalies in acquired sensor data.
  \item Served as data custodian/steward for 10+ TB of oceanographic and atmospheric sensor data; maintained fully documented, reproducible data processing pipelines and version-controlled analysis workflows.
  \item Produced detailed technical reports, visualizations, and safety recommendations from sensor analysis for operational and scientific leadership decision-making.
\end{itemize}

\textbf{Research Physicist~-- Naval Research Laboratory} \hfill 01/2014 -- 06/2016
\vspace*{-2mm}
\begin{itemize}[label={-}, itemsep=-1.0ex]
  \item Designed original Monte Carlo algorithms for probabilistic safety risk quantification and radiological threat detection modeling in Python and C++; evaluated algorithm performance against empirical benchmarks, iterated based on results, and documented findings for reproducibility and peer review -- directly matching the job requirement for 4+ years designing and evaluating data analysis algorithms.
  \item Performed multi-stage predictive model validation; applied statistical methods to identify risk drivers and anomalies in complex sensor measurement data; translated results into technical reports and executive briefings for non-technical government stakeholders.
  \item Results secured \$500K in federal safety R\&D funding. Co-authored peer-reviewed publication (\textit{Journal of Instrumentation}, 2016); operated in classified national security environment under strict safety, governance, and documentation requirements.
\end{itemize}

\textbf{Aerospace Ground Equipment Craftsman~|~Staff Sergeant (E-5)~-- United States Air Force} \hfill 05/2002 -- 02/2010
\vspace*{-2mm}
\begin{itemize}[label={-}, itemsep=-1.0ex]
  \item Maintained and operated safety-critical aerospace ground support equipment for F-16 fighter jet operations worldwide -- direct aviation domain experience with aircraft maintenance safety systems, readiness monitoring, and equipment lifecycle management.
  \item Administered \$2.3M equipment database; led technician teams; reported aircraft safety readiness statistics to base commanders. Honorably discharged.
\end{itemize}

% ML AND ANALYTICS PROJECTS-------------------------------
\vspace*{2mm}
\section{ML and Analytics Projects}
\vspace*{1mm}

\textbf{Boeing Defense~-- Aerial Safety Vision AI} \textit{| Boeing-Sponsored Capstone} \hfill 08/2025 -- 05/2026
\vspace*{-2mm}
\begin{itemize}[label={-}, itemsep=-1.0ex]
  \item Contributed to Boeing Defense-sponsored program designing ML models for aerial drone video analysis -- safety-relevant computer vision applied to aerial imagery. Trained generative models on 263,000+ image datasets via PyTorch DDP on NVIDIA H100/A100 GPU clusters; designed evaluation frameworks (LPIPS, PSNR, T-LPIPS) benchmarked across 10+ real-world flight scenarios.
  \item Maintained fully reproducible training pipeline: version-controlled code (Git), deterministic evaluation protocol (fixed seed 1337), documented checkpointing cadence (200 kimg), and per-tick stats logging -- enabling exact reproduction of any training run. Applied NLP embedding techniques for semantic feature analysis of aerial imagery metadata.
\end{itemize}

\textbf{Sensor Data Analysis~-- LHC Particle Physics~| University of Maryland} \hfill 2013 -- 2015
\vspace*{-2mm}
\begin{itemize}[label={-}, itemsep=-1.0ex]
  \item Designed and implemented gradient-based solvers, efficiency quantification algorithms, and statistical significance methods to extract signal from multi-terabyte particle physics sensor datasets at CERN; evaluated and compared multiple algorithmic approaches against ground truth and published results.
  \item Maintained reproducible analysis workflows using version-controlled code, documented Jupyter notebooks, and shared computing grid environments; co-authored peer-reviewed publication on detector performance analysis.
\end{itemize}

% EDUCATION-------------------------------
\vspace*{2mm}
\section{Education}
\vspace*{1mm}
\textbf{B.S. in Data Analytics}, Washington State University, Everett \hfill 05/2026\\
\textbf{A.A.S. in Computer Science}, Everett Community College \hfill 04/2024\\
GPA: 4.0 (President's List)\\
\textbf{B.S. in Physics}, University of Maryland, College Park \hfill 12/2015

% CERTIFICATIONS-------------------------------
\vspace*{2mm}
\section{Certifications}
\vspace*{1mm}
\textbf{Network Security+}, CompTIA (07/2024)\\
\textbf{Cybersecurity Analyst}, Everett Community College (07/2024)\\
\textbf{Data Analytics}, Washington State University (12/2024)\\
\textbf{Medical Person in Charge}, U.S. Coast Guard Certified

% PUBLICATION-------------------------------
\vspace*{2mm}
\section{Publication}
\vspace*{1mm}
\textit{Journal of Instrumentation} (2016)~-- Liquid Scintillator Tiles for Calorimetry, Large Hadron Collider detector performance analysis

\end{document}